% ====================================================================
% SECTION 04: ARCHITECTURAL INVARIANCE
% ====================================================================

\section{Architectural Invariance of the Duty Cycle}
\label{sec:arch}

\textbf{Mathematical foundation: Robust Optimization.} A Lagrangian structure of
the form $\mathcal{L}_{\mathrm{net}}(\alpha,\eta) =
\eta\,\mathcal{C}_{\mathrm{wave}} + (1-\eta)\,\mathcal{C}_{\mathrm{transport}}$
with $\eta\in[0,1]$ corresponds, in the language of robust control
theory~\cite{ben-tal2009}, to a \emph{Minimax Fractional Excess} problem: one
minimizes, with respect to $\alpha$, the worst-case fractional excess among
incommensurable cost channels. The selection of $\eta$ as a linear weight is the
dual representation of a \emph{worst-case} scenario uncertainty, where Nature
selects the regime (metabolic or pulsatile) that produces the maximum penalty.
The saddle point $(\alpha^*,\eta^*)$ satisfies the cost-balancing condition,
which is the analog of the robust optimality condition for non-fungible
constraints. Hereafter, we use "minimax" as a synonym for this construction,
consistent with the terminology of robust programming.

\chRB{It helps to read this construction as a game. The network chooses the
branching geometry $\alpha$; an adversary---standing in for the uncontrolled
physiological state (rest versus exercise, growth, a changing heart rate)---then
chooses the weight $\eta$ that makes the chosen geometry look as costly as
possible, by tilting the balance toward whichever cost channel that geometry
handles worst. A tree that is excellent for steady transport but reflects pulse
waves can be punished by an adversary who emphasizes the wave channel, and a tree
tuned to suppress reflections can be punished by one who emphasizes transport.
The network's safest move is therefore the geometry at which the adversary's
choice no longer matters---where the two costs are \emph{exactly equal}, so
tilting $\eta$ in either direction changes nothing. That equal-cost geometry is
the saddle point $(\alpha^*,\eta^*)$, and its defining property is \emph{immunity}
to the operating weight: the architecture is optimized not for any single regime
but against the worst case across all of them. This is the precise sense in which
$\alpha^*$ is robust rather than fine-tuned.}

\begin{theorem}[Architectural Invariance of the Biological Duty Cycle]
\label{thm:arch}
Let $(\alpha^*, \eta^*)$ be the unique saddle point of the unified zero-sum
network Lagrangian:
\begin{equation}
    \mathcal{L}_{\mathrm{net}}(\alpha, \eta) =
    \eta\,\mathcal{C}_{\mathrm{wave}}^{\mathrm{net}}(\alpha)
    + (1-\eta)\,\mathcal{C}_{\mathrm{transport}}^{\mathrm{net}}(\alpha).
    \label{eq:lagrangian_net}
\end{equation}
The emergent \chA{marginal-balance weight} $\eta^*$ is an exact invariant of the network's
allometric class $\mathcal{A}(G, p, \alpha_w)$, strictly orthogonal to the
absolute scale of the extensive metabolic parameters $\Lambda\vec{\theta}$.
\end{theorem}

\begin{proof}
The Lagrangian $\mathcal{L}_{\mathrm{net}}(\alpha,\eta)$ is affine (hence both
convex and concave) in $\eta$ and \chA{convex} in $\alpha$ for every fixed $\eta
\in [0,1]$: $C_{\mathrm{transport}}$ is strictly convex, while
$C_{\mathrm{wave}}$ \chA{has $\partial^2 C_{\mathrm{wave}}/\partial\alpha^2 > 0$
throughout the physiological range (verified below), hence is convex there,} with
a unique minimum at $\alpha_w$. The
domains $[\alpha_w, \alpha_t]$ and $[0,1]$ are compact, and
$\mathcal{L}_{\mathrm{net}}$ is continuous. \chA{Because
$\mathcal{L}_{\mathrm{net}}$ is linear in $\eta$, the worst-case cost is attained
at an endpoint,
\begin{equation*}
  \max_{\eta\in[0,1]}\mathcal{L}_{\mathrm{net}}(\alpha,\eta)
  = \max\!\bigl(\mathcal{C}_{\mathrm{wave}}(\alpha),\,
                \mathcal{C}_{\mathrm{transport}}(\alpha)\bigr),
\end{equation*}
the pointwise maximum of two convex functions and hence itself convex in
$\alpha$, so its minimum over the compact interval $[\alpha_w,\alpha_t]$ exists.
Since $\mathcal{C}_{\mathrm{wave}}$ increases away from $\alpha_w$ whereas
$\mathcal{C}_{\mathrm{transport}}$ decreases toward $\alpha_t$ on
$[\alpha_w,\alpha_t]$, and the two costs cross exactly once in the interior, this
minimiser is unique and satisfies the equal-cost condition
$\mathcal{C}_{\mathrm{wave}}(\alpha^*)=\mathcal{C}_{\mathrm{transport}}(\alpha^*)$,
giving the saddle point $(\alpha^*,\eta^*)$. The affine-in-$\eta$,
convex-in-$\alpha$ structure thus reduces existence and uniqueness of the saddle
to elementary convex analysis, with no appeal to the general
quasiconvex--quasiconcave minimax theorem.} Detailed
numerical verification (Supplemental Material, Table S1) confirms that the
\chA{scalar} Hessian $\partial^2 \mathcal{C}_{\mathrm{wave}}/\partial\alpha^2 > 0$ throughout
the physiological range $\alpha \in [2.0, 3.5]$ and $\eta \in [0,1]$,
\chA{securing the convexity on which this argument rests.}

Uniqueness follows from the \textbf{strict convexity} of
$\mathcal{C}_{\mathrm{transport}}$ and the \textbf{strict monotonicity} of the
marginal penalties: because $\partial
\mathcal{C}_{\mathrm{transport}}^{\mathrm{net}}/\partial \alpha$ decreases
monotonically from zero (becomes more negative) as $\alpha$ moves below
$\alpha_t$, and $\partial \mathcal{C}_{\mathrm{wave}}^{\mathrm{net}}/\partial
\alpha$ increases monotonically from zero (becomes more positive) as $\alpha$
moves above $\alpha_w$, their intersection at any fixed $\eta$ is unique. The
saddle-point condition $\partial\mathcal{L}_{\mathrm{net}}/\partial\alpha = 0$
yields:
\begin{equation}
    \eta^* = \left( 1 + \frac{\partial \mathcal{C}_{\mathrm{wave}}^{\mathrm{net}}
    (\alpha^*) / \partial \alpha}
    {|\partial \mathcal{C}_{\mathrm{transport}}^{\mathrm{net}}(\alpha^*)
    / \partial \alpha|} \right)^{-1}.
    \label{eq:eta_star}
\end{equation}
\chA{Because $\mathcal{L}_{\mathrm{net}}$ is affine in $\eta$, an interior saddle
with $\eta^*\in(0,1)$ can occur only when the maximiser is indifferent between the
two endpoint costs; equivalently,
\begin{equation*}
  \frac{\partial\mathcal{L}_{\mathrm{net}}}{\partial\eta}
  = \mathcal{C}_{\mathrm{wave}}^{\mathrm{net}}(\alpha^*)
  - \mathcal{C}_{\mathrm{transport}}^{\mathrm{net}}(\alpha^*) = 0,
\end{equation*}
which is the equal-cost condition.} \chRB{This is the
formal counterpart of the game intuition above: at the saddle the adversary is
indifferent to $\eta$ precisely because the two costs coincide.}
This equivalence
ensures that the emergent \chA{marginal-balance weight} $\eta^*$ effectively balances
the two incommensurable penalties. The numerator derives from the acoustic impedance
mismatch at bifurcations and is a strictly geometric function of $(G, N,
\alpha_w, \alpha^*)$ containing no metabolic parameters.

By Theorem~\ref{thm:gauge}, $\mathcal{C}_{\mathrm{transport}}^{\mathrm{net}}$ is
a degree-zero homogeneous function of $\vec{\theta}$. Its derivative with
respect to $\alpha$ is therefore identically gauge-invariant:
\begin{equation}
    \frac{\partial\mathcal{C}_{\mathrm{transport}}^{\mathrm{net}}}
    {\partial\alpha}(\alpha; \Lambda\vec{\theta})
    = \frac{\partial\mathcal{C}_{\mathrm{transport}}^{\mathrm{net}}}
    {\partial\alpha}(\alpha; \vec{\theta}).
\end{equation}
In Eq.~\eqref{eq:eta_star} the watts cancel identically in the gradient ratio.
The \chA{marginal-balance weight} $\eta^*$ is therefore a pure structural invariant,
determined solely by $(G, p, \alpha_w)$ through the geometric evaluation of both
derivatives at $\alpha^*$.
\end{proof}

\begin{remark}
This provides a first-principles explanation for the empirical observation that
vascular branching exponents are conserved across developmental stages within a
species, as suggested by the generation-invariant morphometry
of~\cite{kassab1993} and the developmental convergence reported
in~\cite{forjaz2026}.
\end{remark}

\begin{remark}[Theoretical Bounds on the Static Transport Attractor]
Paper~I~\cite{paperI} establishes that the static transport attractor $\alpha_t$
must lie within strict theoretical bounds derived from three-term metabolic cost
structure: $(5+p)/2 \approx 2.89 < \alpha_t < 3$ (Theorem~3.2), independent of
flow asymmetry. Empirical values---ranging from $\alpha^* \approx
\VarAlphaRatObs$ in rat pulmonary arteries ($M=\VarMassRat\,$g)~\cite{jiang1994}
to $\alpha^* \approx \VarAlphaStar$ in large mammals~\cite{kassab1993}---lie
below this static bound precisely because pulsatile dynamics contribute an
additional wave-reflection penalty missing from purely viscous optimization,
demonstrating that the minimax solution is the inevitable physical attractor
within these constraints.
\end{remark}

\begin{remark}[Convexity Verification: Physical and Numerical]
\label{rem:quasiconvex}
\chA{Existence and uniqueness of the saddle in Theorem~\ref{thm:arch} require only
that $\mathcal{C}_{\mathrm{wave}}(\alpha)$ be convex in $\alpha$ for each fixed
$\eta$. We establish this} through the following physical and numerical arguments:

\textbf{Physical argument.} The wave-reflection cost measures impedance mismatch
at vascular junctions. The wave attractor exponent $\alpha_w = 2$ emerges from
global reflection minimization (Paper~II~\cite{paperII}): at each bifurcation,
the power reflection coefficient $R^2(\alpha) =
[(\gamma(\alpha)-1)/(\gamma(\alpha)+1)]^2$ with $\gamma(\alpha) =
N^{1-2/\alpha}$ vanishes uniquely at $\alpha = 2$, selecting this value as the
dynamically stable attractor over the network hierarchy. Away from this
geometric impedance-matching condition ($\alpha = \alpha_w = 2$), the reflection
coefficient $|\Gamma(\alpha)|^2$ grows \emph{monotonically}: larger $\alpha$
(narrower daughters) increases mismatch in one direction, smaller $\alpha$
(wider daughters) increases it in the opposite direction. This monotonic growth
ensures no local minima can exist. The cumulative reflection penalty
$\mathcal{C}_{\mathrm{wave}}(\alpha) \propto \sum_g |\Gamma_g(\alpha)|^2$
inherits this unimodal structure. Crucially, viscous damping (attenuation factor
$e^{-2\kappa L_g}$ in the wave propagation) and spectral averaging over cardiac
harmonics eliminate coherent interference effects that might otherwise introduce
secondary extrema via resonances.

\textbf{Numerical verification.} Detailed numerical evaluation of
$\mathcal{C}_{\mathrm{wave}}(\alpha)$ for the physiological parameter range
$\alpha \in [2.0, 3.5]$ and $\eta \in [0,1]$ confirms strict unimodality
(Supplemental Material, Table S1). No spurious local minima or inflection points
are observed. The transport cost $\mathcal{C}_{\mathrm{transport}}(\alpha)$ is
\emph{strictly} convex (proven analytically in Paper~II~\cite{paperII} via
power-law form in viscous dissipation). \chA{Since $\mathcal{C}_{\mathrm{wave}}$
is convex over the physiological range and $\mathcal{C}_{\mathrm{transport}}$ is
strictly convex, their non-negative combination
$\mathcal{L}_{\mathrm{net}}(\alpha,\eta)$ is convex in $\alpha$, so the
convex-envelope argument of Theorem~\ref{thm:arch} yields a unique saddle
point.} The
minimax equal-cost condition $f(\alpha^*) = g(\alpha^*)$ is verified without
free parameters (Paper~III~\cite{paperIII}): for the porcine coronary tree,
solving this balance yields $\alpha^* = 2.77$, consistent with the morphometric
average $\alpha_{\exp} = 2.70 \pm 0.20$, confirming that the network-level
Lagrangian balances measurable first-principles costs.
\end{remark}

\begin{remark}[Allometric vs.\ angiogenic invariance]
The invariance of $\eta^*$ established in Theorem~\ref{thm:arch} is an
\emph{allometric} invariance: it holds under rescaling of the metabolic
parameters $\vec{\theta}$ at fixed network topology $(G, N, \alpha_w, p)$. It
does not apply to \emph{angiogenic} changes in $G$ (addition or pruning of
vascular generations), which alter $\kappa_{\mathrm{eff}}(G)$ and thereby shift
$\alpha^*$. The two types of invariance are physically and mathematically
distinct; only the former is established here.
\end{remark}

\begin{remark}[Physical status of $\Phi_{\mathrm{opt}}$]
The reference state $\Phi^{\mathrm{net}}(\alpha_{\mathrm{opt}})$ is the cost
attained when every vessel independently sits at its locally optimal radius
$r^*(Q_g)$---a configuration achievable only if the network is allowed to
violate the global self-similar constraint $r_g = r_0 N^{-g/\alpha}$. It is
therefore a mathematical lower bound (a gauge baseline), not a physiologically
realizable state of the intact tree. The fractional excess
$\mathcal{C}^{\mathrm{net}}_{\mathrm{transport}} = (\Phi^{\mathrm{net}} -
\Phi_{\mathrm{opt}})/\Phi_{\mathrm{opt}}$ measures the architectural cost of
imposing global self-similarity, not a deviation from a physically accessible
optimum.
\end{remark}

\begin{remark}[The Static Attractor]
As demonstrated in Paper~I~\cite{paperI}, the explicit evaluation of the static
transport optimum $\alpha_t$---with the inclusion of the metabolic cost of the
vascular wall ($\propto r^{1+p}$)---strictly breaks the universality of Murray's
law ($\alpha=3$) and fixes the purely static attractor at $\alpha_t \approx
\VarAlphaTLow-\VarAlphaTHigh$ for mammalian coronary networks. This value
constitutes the theoretical starting point of the present framework: the
residual gap from $\alpha_t \approx \VarAlphaTLow$ to the empirical $\alpha
\approx \VarAlphaStar$ is a mathematical necessity that requires the inclusion
of the dynamic, pulsatile minimax mechanism developed here.
\end{remark}
